% GENERATED FILE. Edit canonical lessons, then run npm run generate:books.
% canonical-source-hash: fnv1a64:36a28da79c3de0ec
% canonical-lessons: ES-C552-inmigrante, ES-C552-exilio, ES-C552-refugiado, ES-C552-frontera, ES-C552-independizarse

\chapter{Lejos de casa}
\label{ch:es-lejos-de-casa}
\hlchaptermodality{552}

\begin{chapteropening}
\textbf{By the end of this chapter:} \emph{I can talk about people who left their country: immigrants, refugees, exile, borders and leaving home.}

Complete the last lesson of chapter 552: I can talk about people who left their country: immigrants, refugees, exile, borders and leaving home.
\end{chapteropening}

\section[el inmigrante]{el inmigrante — the immigrant}

\label{lesson:ES-C552-inmigrante}



\begin{quote}
Before the new one: say the Spanish for to emigrate, then the Spanish for to adapt.
\end{quote}

\subsection*{You'll want to know: el inmigrante}
\textbf{el inmigrante} — \textquotedblleft{}the immigrant\textquotedblright{}.

\begin{quote}
Muchos inmigrantes trabajan en el campo.
\end{quote}

\begin{quote}
\emph{Many immigrants work on the land.}
\end{quote}

The same form for a woman: \emph{la inmigrante}.

\begin{cousinweb}
\emph{inmigrante} is built on \emph{inmigrar}, a learned word from Latin \emph{immigrāre} 'to move into' (Lewis \& Short), built on \emph{in} and \emph{migrāre} 'to move'.
\end{cousinweb}

\begin{grammarlens}[title={What you've built}]
One more word for telling, explaining and describing at B1.
\end{grammarlens}

\subsection*{Your turn}
\begin{itemize}
\raggedright
  \item \emph{Say it:} \emph{el inmigrante}
  \item \emph{Say it:} \emph{el inmigrante}, once more
  \item \emph{Say it:} a sentence of your own with \emph{el inmigrante}
  \item \emph{Recall:} say the Spanish for to emigrate, then the Spanish for to adapt, then say \emph{el inmigrante} again
\end{itemize}

\begin{tcolorbox}[breakable,colback=teal!4,colframe=teal!35!black,title={Before you move on}]
What does el inmigrante mean? (\textquotedblleft{}The immigrant\textquotedblright{}.) Say it once more.
\end{tcolorbox}
\section[el exilio]{el exilio — the exile}

\label{lesson:ES-C552-exilio}



\begin{quote}
Before the new one: say the Spanish for to adapt, then the Spanish for the immigrant.
\end{quote}

\subsection*{You'll want to know: el exilio}
\textbf{el exilio} — \textquotedblleft{}the exile\textquotedblright{}.

\begin{quote}
Vivió veinte años en el exilio.
\end{quote}

\begin{quote}
\emph{He lived in exile for twenty years.}
\end{quote}

\emph{Exiliarse} is to go into exile.

\begin{cousinweb}
\emph{exilio} is a learned borrowing from Latin \emph{exilium} 'banishment' (EtymDB).
\end{cousinweb}

\begin{grammarlens}[title={What you've built}]
One more word for telling, explaining and describing at B1.
\end{grammarlens}

\subsection*{Your turn}
\begin{itemize}
\raggedright
  \item \emph{Say it:} \emph{el exilio}
  \item \emph{Say it:} \emph{el exilio}, once more
  \item \emph{Say it:} a sentence of your own with \emph{el exilio}
  \item \emph{Recall:} say the Spanish for to adapt, then the Spanish for the immigrant, then say \emph{el exilio} again
\end{itemize}

\begin{tcolorbox}[breakable,colback=teal!4,colframe=teal!35!black,title={Before you move on}]
What does el exilio mean? (\textquotedblleft{}The exile\textquotedblright{}.) Say it once more.
\end{tcolorbox}
\section[el refugiado]{el refugiado — the refugee}

\label{lesson:ES-C552-refugiado}



\begin{quote}
Before the new one: say the Spanish for the immigrant, then the Spanish for the exile.
\end{quote}

\subsection*{You'll want to know: el refugiado}
\textbf{el refugiado} — \textquotedblleft{}the refugee\textquotedblright{}.

\begin{quote}
El país acogió a miles de refugiados.
\end{quote}

\begin{quote}
\emph{The country took in thousands of refugees.}
\end{quote}

\begin{cousinweb}
\emph{refugiado} is built on \emph{refugiarse} 'to take shelter', from \emph{refugio}, which comes from Latin \emph{refugium} 'a taking refuge' (Lewis \& Short), built on \emph{fugere} 'to flee'.
\end{cousinweb}

\begin{grammarlens}[title={What you've built}]
One more word for telling, explaining and describing at B1.
\end{grammarlens}

\subsection*{Your turn}
\begin{itemize}
\raggedright
  \item \emph{Say it:} \emph{el refugiado}
  \item \emph{Say it:} \emph{el refugiado}, once more
  \item \emph{Say it:} a sentence of your own with \emph{el refugiado}
  \item \emph{Recall:} say the Spanish for the immigrant, then the Spanish for the exile, then say \emph{el refugiado} again
\end{itemize}

\begin{tcolorbox}[breakable,colback=teal!4,colframe=teal!35!black,title={Before you move on}]
What does el refugiado mean? (\textquotedblleft{}The refugee\textquotedblright{}.) Say it once more.
\end{tcolorbox}
\section[la frontera]{la frontera — the border, the frontier}

\label{lesson:ES-C552-frontera}



\begin{quote}
Before the new one: say the Spanish for the exile, then the Spanish for the refugee.
\end{quote}

\subsection*{You'll want to know: la frontera}
\textbf{la frontera} — \textquotedblleft{}the border, the frontier\textquotedblright{}.

\begin{quote}
Cruzaron la frontera de noche.
\end{quote}

\begin{quote}
\emph{They crossed the border at night.}
\end{quote}

\begin{cousinweb}
\emph{frontera} is usually traced to an old word for 'facing, opposite', from Latin \emph{frōns}, \emph{frontem} 'forehead, front' (EtymDB). The border is the line where two countries face each other.
\end{cousinweb}

\begin{grammarlens}[title={What you've built}]
One more word for telling, explaining and describing at B1.
\end{grammarlens}

\subsection*{Your turn}
\begin{itemize}
\raggedright
  \item \emph{Say it:} \emph{la frontera}
  \item \emph{Say it:} \emph{la frontera}, once more
  \item \emph{Say it:} a sentence of your own with \emph{la frontera}
  \item \emph{Recall:} say the Spanish for the exile, then the Spanish for the refugee, then say \emph{la frontera} again
\end{itemize}

\begin{tcolorbox}[breakable,colback=teal!4,colframe=teal!35!black,title={Before you move on}]
What does la frontera mean? (\textquotedblleft{}The border, the frontier\textquotedblright{}.) Say it once more.
\end{tcolorbox}
\section[independizarse]{independizarse — to become independent; to leave home}

\label{lesson:ES-C552-independizarse}



\begin{quote}
Before the new one: say the Spanish for the refugee, then the Spanish for the border.
\end{quote}

\subsection*{You'll want to know: independizarse}
\textbf{independizarse} — \textquotedblleft{}to become independent; to leave home\textquotedblright{}.

\begin{quote}
Me independicé a los veintidós años.
\end{quote}

\begin{quote}
\emph{I left home at twenty-two.}
\end{quote}

For a young person it means moving out of the parents' house. It takes \emph{c} before \emph{e}: \emph{me independicé}.

\begin{cousinweb}
\emph{independizarse} is built on \emph{independiente} with the ending \emph{-izar}. The root is \emph{depender}, from Latin \emph{dēpendēre} 'to hang down from', built on \emph{pendēre} 'to hang' (Lewis \& Short).
\end{cousinweb}

\begin{grammarlens}[title={What you've built}]
One more word for telling, explaining and describing at B1.
\end{grammarlens}

\subsection*{Your turn}
\begin{itemize}
\raggedright
  \item \emph{Say it:} \emph{independizarse}
  \item \emph{Say it:} \emph{independizarse}, once more
  \item \emph{Say it:} a sentence of your own with \emph{independizarse}
  \item \emph{Recall:} say the Spanish for the refugee, then the Spanish for the border, then say \emph{independizarse} again
\end{itemize}

\begin{tcolorbox}[breakable,colback=teal!4,colframe=teal!35!black,title={Before you move on}]
What does independizarse mean? (\textquotedblleft{}To become independent; to leave home\textquotedblright{}.) Say it once more.
\end{tcolorbox}
